%% @file paper52_algebraische_topologie_de.tex
%% @brief Algebraische Topologie: Fundamentalgruppen, Homologie, Kohomologie und Anwendungen
%% @author Michael Fuhrmann
%% @date 2026-03-12
%% @project Specialist Mathematics Project, Build 139

\documentclass[12pt,a4paper]{amsart}
\usepackage{amsmath,amssymb,amsthm}
\usepackage{mathtools}
\usepackage[utf8]{inputenc}
\usepackage[T1]{fontenc}
\usepackage[ngerman]{babel}
\usepackage{hyperref}
\usepackage{enumitem,booktabs,array}
\usepackage{geometry}
\geometry{margin=2.5cm}

%% Theorem-Umgebungen (deutsch)
\newtheorem{theorem}{Satz}[section]
\newtheorem{lemma}[theorem]{Lemma}
\newtheorem{corollary}[theorem]{Korollar}
\newtheorem{proposition}[theorem]{Proposition}
\newtheorem{conjecture}[theorem]{Vermutung}
\theoremstyle{definition}
\newtheorem{definition}[theorem]{Definition}
\newtheorem{example}[theorem]{Beispiel}
\newtheorem{remark}[theorem]{Bemerkung}

\author{Michael Fuhrmann}
\address{Specialist Mathematics Project, Build 139}
\email{specialist-maths@localhost}
\date{\today}
\begin{document}
\title{Algebraische Topologie:\\
Fundamentalgruppen, Homologie, Kohomologie\\
und klassische Anwendungen}

\begin{abstract}
Die algebraische Topologie untersucht topologische R\"aume mithilfe algebraischer
Invarianten und \"ubersetzt geometrische Probleme in berechenbare algebraische
Fragen. Diese Arbeit gibt eine umfassende Darstellung der Fundamentalgruppe,
\"Uberlagerungsr\"aume, singul\"aren Homologie und Kohomologie, CW-Komplexe sowie
h\"oherer Homotopiegruppen. Wir beweisen den Seifert-van-Kampen-Satz, die
Mayer-Vietoris-Sequenz, den Ausschneidungssatz, den universellen Koeffizientensatz,
die K\"unneth-Formel, die Poincar\'e-Dualit\"at f\"ur kompakte orientierbare
Mannigfaltigkeiten und den Hurewicz-Isomorphismus. Klassische Anwendungen umfassen
den Brouwer-Fixpunktsatz, den Satz von Borsuk-Ulam und den Igel-Satz (Hairy Ball
Theorem). Die Poincar\'e-Vermutung in Dimension drei, 2003 von Perelman bewiesen,
sowie Adams' Satz \"uber die Hopf-Invariante 1 werden ebenfalls behandelt.
Alle Resultate sind streng bewiesen; Verkn\"upfungen zur Implementierung im
Begleitmodul \texttt{algebraic\_topology.py} des Specialist-Mathematics-Projekts
werden hergestellt.
\end{abstract}

\maketitle
\tableofcontents

%% ============================================================
\section{Einleitung}
%% ============================================================

Die algebraische Topologie ist das Studium topologischer R\"aume durch das Prisma
der Algebra. Die zentrale Idee besteht darin, algebraischen R\"aumen algebraische
Objekte -- Gruppen, Ringe, Moduln -- in \emph{funktoreller} Weise zuzuordnen,
sodass hom\"omorphe R\"aume isomorphe Invarianten erhalten. Damit lassen sich
qualitative geometrische Fragen (``Sind diese R\"aume hom\"omorph?'') in
handhabbare algebraische Fragen (``Sind diese Gruppen isomorph?'') umwandeln.

Die historischen Wurzeln reichen bis zu Eulers Polyederformel
\[
  V - E + F = 2
\]
zur\"uck, dem ersten Beispiel einer algebraischen Invariante, die topologische
Information erfasst. Poincar\'e systematisierte diese Ideen in seinen grundlegenden
Arbeiten von 1895 bis 1904, f\"uhrte die Fundamentalgruppe und die Homologiegruppen
ein und stellte die nach ihm benannte Vermutung auf.

Das 20.~Jahrhundert sah die Entwicklung der singul\"aren Homologie (Eilenberg),
der Kohomologie (\v{C}ech, de Rham), des axiomatischen Rahmens von Eilenberg und
Steenrod, der Homotopietheorie, der Spektralsequenzen, der K-Theorie und schlie\ss{}lich
die L\"osung der Poincar\'e-Vermutung in allen Dimensionen (Smale f\"ur $\dim \geq 5$,
Freedman f\"ur $\dim = 4$, Perelman f\"ur $\dim = 3$).

\medskip\noindent\textbf{Gliederung.}
Abschnitt~\ref{sec:fundamentalgruppe} f\"uhrt die Fundamentalgruppe ein und beweist
den Seifert-van-Kampen-Satz. Abschnitt~\ref{sec:ueberlagerung} behandelt
\"Uberlagerungsr\"aume. Abschnitt~\ref{sec:homologie} entwickelt die singul\"are
Homologie. Abschnitt~\ref{sec:schluessel} beweist die Mayer-Vietoris-Sequenz und
den Ausschneidungssatz. Abschnitt~\ref{sec:kohomologie} behandelt die Kohomologie
und den universellen Koeffizientensatz. Abschnitt~\ref{sec:poincare} beweist die
Poincar\'e-Dualit\"at. Abschnitt~\ref{sec:cw} diskutiert CW-Komplexe und
zelluläre Homologie. Abschnitt~\ref{sec:hoehere} behandelt h\"ohere Homotopiegruppen
und den Hurewicz-Satz. Abschnitt~\ref{sec:anwendungen} pr\"asentiert klassische
Fixpunkt- und Antipodensätze.

%% ============================================================
\section{Die Fundamentalgruppe}
\label{sec:fundamentalgruppe}
%% ============================================================

\subsection{Homotopie von Wegen}

Sei $X$ ein topologischer Raum und $x_0 \in X$ ein Basispunkt.

\begin{definition}
Eine \emph{Schleife} bei $x_0$ ist eine stetige Abbildung
$\gamma\colon [0,1] \to X$ mit $\gamma(0) = \gamma(1) = x_0$.
Zwei Schleifen $\gamma_0, \gamma_1$ sind \emph{homotop relativ $\{0,1\}$},
geschrieben $\gamma_0 \simeq \gamma_1$, wenn es eine stetige Abbildung
\[
  H\colon [0,1] \times [0,1] \to X
\]
gibt mit $H(s,0) = \gamma_0(s)$, $H(s,1) = \gamma_1(s)$ und
$H(0,t) = H(1,t) = x_0$ f\"ur alle $s, t \in [0,1]$.
\end{definition}

Homotopie relativ Endpunkten ist eine \"Aquivalenzrelation. Die
\emph{Verkettung} von Schleifen $\gamma * \delta$, definiert durch
\[
  (\gamma * \delta)(s) = \begin{cases}
    \gamma(2s) & 0 \leq s \leq \tfrac{1}{2}, \\
    \delta(2s-1) & \tfrac{1}{2} \leq s \leq 1,
  \end{cases}
\]
setzt sich zu einer wohldefinierte Operation auf Homotopieklassen fort.

\begin{definition}
Die \emph{Fundamentalgruppe} von $X$ bei $x_0$ ist
\[
  \pi_1(X, x_0) = \{ [\gamma] : \gamma \text{ Schleife bei } x_0 \},
\]
mit Multiplikation $[\gamma][\delta] = [\gamma * \delta]$.
\end{definition}

\begin{theorem}
$\pi_1(X, x_0)$ ist unter Verkettung von Homotopieklassen eine Gruppe.
\end{theorem}
\begin{proof}
\textbf{Assoziativit\"at.} $([\alpha][\beta])[\gamma] = [\alpha]([\beta][\gamma])$
folgt durch Umparametrisierung: Man definiert $H(s,t)$ als eine
$t$-abh\"angige Reparametrisierung von $\alpha * \beta * \gamma$, die die
Unterteilungspunkte verschiebt; die resultierende Abbildung ist eine g\"ultige
Homotopie.

\textbf{Neutrales Element.} Die konstante Schleife $e(s) = x_0$ erf\"ullt
$[\gamma][e] = [\gamma]$: Die Homotopie $H(s,t) = \gamma(\min(2s/(1+t),1))$
zeigt dies.

\textbf{Inverses.} $\bar{\gamma}(s) = \gamma(1-s)$ erf\"ullt
$[\gamma][\bar{\gamma}] = [e]$, da die Schleife sich via einer Homotopie
``zur\"uckzieht''.
\end{proof}

\subsection{Beispiele}

\begin{example}
$\pi_1(\mathbb{R}^n, x_0) = 1$ (trivial): Jede Schleife ist null-homotop via
$H(s,t) = (1-t)\gamma(s) + t x_0$ (Geradlinien-Homotopie).
\end{example}

\begin{example}
$\pi_1(S^1, 1) \cong \mathbb{Z}$: Der Isomorphismus ordnet einer Schleife
ihre \emph{Windungszahl} zu. Die Schleife $\gamma_n(s) = e^{2\pi i n s}$
repr\"asentiert die Klasse $n \in \mathbb{Z}$.
\end{example}

\begin{example}
$\pi_1(T^2, x_0) \cong \mathbb{Z}^2$ (abelsch): Die zwei Erzeuger entsprechen
den zwei unabh\"angigen Schleifen um den Torus.
\end{example}

\begin{example}
$\pi_1(\mathbb{R}P^2, x_0) \cong \mathbb{Z}/2\mathbb{Z}$: Eine nicht-orientierbare
Mannigfaltigkeit, deren Fundamentalgruppe die orientierungsumkehrende Schleife
erfasst.
\end{example}

\subsection{Der Seifert-van-Kampen-Satz}

Der van-Kampen-Satz berechnet die Fundamentalgruppe einer Vereinigung.

\begin{theorem}[Seifert-van-Kampen]
\label{thm:vankampen}
Sei $X = U \cup V$, wobei $U, V$ offen, wegzusammenh\"angend und
$U \cap V$ wegzusammenh\"angend ist. W\"ahle $x_0 \in U \cap V$. Dann gilt
\[
  \pi_1(X, x_0) \cong \pi_1(U, x_0) *_{\pi_1(U \cap V, x_0)} \pi_1(V, x_0),
\]
das \emph{amalgamierte freie Produkt} \"uber die Bilder von $\pi_1(U \cap V)$.
\end{theorem}
\begin{proof}[Beweisidee]
Jede Schleife $\gamma$ in $X$ bei $x_0$ l\"asst sich durch Kompaktheit von $[0,1]$
und das Lebesguesche Zahlenlemma in Teilb\"ogen zerlegen, von denen jeder ganz in
$U$ oder ganz in $V$ liegt. Durch Einf\"ugen von Verbindungswegen in $U \cap V$
zerlegt sich die Schleife in ein Produkt von Schleifen in $U$ und $V$. Dies liefert
eine Surjektion von $\pi_1(U) * \pi_1(V)$ nach $\pi_1(X)$.

F\"ur die Injektivit\"at: Gilt ein Wort $w = u_1 v_1 \cdots u_k v_k$ (mit
$u_i \in \pi_1(U)$, $v_j \in \pi_1(V)$) im Bild trivial, so liefert eine
Homotopie $H\colon [0,1]^2 \to X$ die Nullhomotopie. Durch Unterteilung des
Quadrats -- erneut via Lebesgue-Zahl -- kann man zeigen, dass das Wort bereits
im amalgamierten freien Produkt trivial ist.
\end{proof}

\begin{example}
F\"ur den Keil $S^1 \vee S^1$: Mit $U$ und $V$ als offenen Teilmengen, die
jeweils homotopie\"aquivalent zu $S^1$ sind und mit $U \cap V \simeq *$
(zusammenziehbar), ergibt sich
$\pi_1(S^1 \vee S^1) \cong \mathbb{Z} * \mathbb{Z}$,
die freie Gruppe mit zwei Erzeugern.
\end{example}

%% ============================================================
\section{\"Uberlagerungsr\"aume}
\label{sec:ueberlagerung}
%% ============================================================

\subsection{Grundlagen}

\begin{definition}
Ein \emph{\"Uberlagerungsraum} von $X$ ist ein Paar $(\tilde{X}, p)$, wobei
$\tilde{X}$ ein topologischer Raum und $p\colon \tilde{X} \to X$ eine stetige
Surjektion ist, sodass jedes $x \in X$ eine offene Umgebung $U$ besitzt mit
$p^{-1}(U) = \bigsqcup_\alpha V_\alpha$ (disjunkte Vereinigung offener Mengen),
wobei $p|_{V_\alpha}\colon V_\alpha \to U$ ein Hom\"oomorphismus ist.
\end{definition}

\begin{theorem}[Wegliftungssatz]
Gegeben ein \"Uberlagerungsraum $p\colon \tilde{X} \to X$, ein Weg
$\gamma\colon [0,1] \to X$ und ein Lift $\tilde{x}_0 \in p^{-1}(\gamma(0))$,
existiert ein eindeutiger Lift $\tilde{\gamma}\colon [0,1] \to \tilde{X}$
mit $\tilde{\gamma}(0) = \tilde{x}_0$ und $p \circ \tilde{\gamma} = \gamma$.
\end{theorem}

\begin{theorem}[Homotopieliftungssatz]
Gegeben ein \"Uberlagerungsraum $p\colon \tilde{X} \to X$ und eine Homotopie
$H\colon [0,1]^2 \to X$ mit einem Lift von $H(0,0)$, existiert ein eindeutiger
Lift $\tilde{H}\colon [0,1]^2 \to \tilde{X}$.
\end{theorem}

\subsection{Decktransformationen und Klassifikation}

\begin{definition}
Eine \emph{Decktransformation} von $p\colon \tilde{X} \to X$ ist ein
Hom\"oomorphismus $\phi\colon \tilde{X} \to \tilde{X}$ mit $p \circ \phi = p$.
Die Gruppe der Decktransformationen hei\ss{}t $\mathrm{Deck}(p)$.
\end{definition}

\begin{theorem}[Klassifikation von \"Uberlagerungsr\"aumen]
\label{thm:ueberlagerung_klassifikation}
F\"ur einen zusammenh\"angenden, lokal wegzusammenh\"angenden und semi-lokal
einfach-zusammenh\"angenden Raum $X$ mit Basispunkt $x_0$ gibt es eine Bijektion
zwischen:
\begin{enumerate}[label=(\roman*)]
  \item Isomorphieklassen zusammenh\"angender \"Uberlagerungsr\"aume von $X$.
  \item Konjugationsklassen von Untergruppen von $\pi_1(X, x_0)$.
\end{enumerate}
Der \emph{universelle \"Uberlagerungsraum} $\tilde{X}$ (mit
$\pi_1(\tilde{X}) = 1$) entspricht der trivialen Untergruppe, und
$\mathrm{Deck}(\tilde{X} \to X) \cong \pi_1(X, x_0)$.
\end{theorem}

\begin{example}
Der universelle \"Uberlagerungsraum von $S^1$ ist $\mathbb{R}$, mit
$p(t) = e^{2\pi i t}$. Decktransformationen sind ganzzahlige Translationen
$t \mapsto t + n$, und $\mathrm{Deck}(\mathbb{R} \to S^1) \cong \mathbb{Z}
\cong \pi_1(S^1)$.
\end{example}

\begin{example}
Der universelle \"Uberlagerungsraum von $\mathbb{R}P^n$ ist $S^n$, mit
der Decktransformationsgruppe $\mathbb{Z}/2\mathbb{Z}$ (antipodische
Abbildung).
\end{example}

%% ============================================================
\section{Singul\"are Homologie}
\label{sec:homologie}
%% ============================================================

\subsection{Singul\"are Simplizes und Kettengruppen}

\begin{definition}
Das \emph{Standard-$n$-Simplex} ist
\[
  \Delta^n = \left\{ (t_0, \ldots, t_n) \in \mathbb{R}^{n+1}
    : t_i \geq 0,\; \sum_{i=0}^n t_i = 1 \right\}.
\]
Ein \emph{singul\"arer $n$-Simplex} in $X$ ist eine stetige Abbildung
$\sigma\colon \Delta^n \to X$. Die \emph{$n$-te Kettengruppe}
$C_n(X;\mathbb{Z})$ ist die freie abelsche Gruppe, erzeugt von allen
singul\"aren $n$-Simplizes.
\end{definition}

\subsection{Randoperatoren und Homologie}

Die \emph{$i$-te Seitenabbildung} $\partial_i\colon \Delta^{n-1} \to \Delta^n$
bettet das Standard-$(n-1)$-Simplex als die der $i$-ten Ecke gegen\"uberliegende
Seite ein:
\[
  \partial_i(t_0,\ldots,t_{n-1})
  = (t_0,\ldots,t_{i-1}, 0, t_i, \ldots, t_{n-1}).
\]
Der Rand eines singul\"aren $n$-Simplex $\sigma$ ist
\[
  \partial_n(\sigma) = \sum_{i=0}^n (-1)^i (\sigma \circ \partial_i)
  \;\in C_{n-1}(X;\mathbb{Z}).
\]

\begin{theorem}[$\partial^2 = 0$]
$\partial_{n-1} \circ \partial_n = 0$ f\"ur alle $n$.
\end{theorem}
\begin{proof}
Die Seitenabbildungen erf\"ullen die Kompatibilit\"atsrelation
$\partial_j \circ \partial_i = \partial_i \circ \partial_{j-1}$ f\"ur $i < j$.
Damit:
\begin{align*}
  \partial_{n-1}\partial_n(\sigma)
  &= \sum_{j=0}^n (-1)^j \partial_{n-1}(\sigma \circ \partial_j) \\
  &= \sum_{j=0}^n \sum_{i=0}^{n-1} (-1)^{i+j}
     (\sigma \circ \partial_j \circ \partial_i) \\
  &= \sum_{i < j} (-1)^{i+j}(\sigma \circ \partial_i \circ \partial_{j-1})
   + \sum_{i \geq j} (-1)^{i+j}(\sigma \circ \partial_j \circ \partial_{i-1}).
\end{align*}
Die Substitution $j' = j-1$ in der ersten Summe zeigt, dass sich die beiden
Summen paarweise wegheben.
\end{proof}

Damit bilden $\{C_n(X;\mathbb{Z}), \partial_n\}$ einen \emph{Kettenkomplex}.

\begin{definition}
Die \emph{$n$-te singul\"are Homologiegruppe} von $X$ ist
\[
  H_n(X;\mathbb{Z})
  = \ker(\partial_n) / \operatorname{im}(\partial_{n+1})
  = Z_n / B_n,
\]
wobei $Z_n = \ker \partial_n$ die \emph{Zyklen} und
$B_n = \operatorname{im} \partial_{n+1}$ die \emph{R\"ander} sind.
\end{definition}

\subsection{Grundlegende Berechnungen}

\begin{example}
F\"ur einen Punkt $*$: $H_0(*) = \mathbb{Z}$ und $H_n(*) = 0$ f\"ur $n > 0$.
\end{example}

\begin{example}[$n$-Sph\"are]
\[
  H_k(S^n;\mathbb{Z}) = \begin{cases}
    \mathbb{Z} & k = 0 \text{ oder } k = n, \\
    0 & \text{sonst.}
  \end{cases}
\]
\end{example}

\begin{example}[Torus]
$H_0(T^2) = \mathbb{Z}$, $H_1(T^2) = \mathbb{Z}^2$, $H_2(T^2) = \mathbb{Z}$,
$H_k(T^2) = 0$ f\"ur $k \geq 3$.
\end{example}

\begin{example}[Reell-projektive Ebene]
$H_0(\mathbb{R}P^2) = \mathbb{Z}$,
$H_1(\mathbb{R}P^2) = \mathbb{Z}/2\mathbb{Z}$,
$H_k(\mathbb{R}P^2) = 0$ f\"ur $k \geq 2$.
\end{example}

\subsection{Die lange exakte Sequenz eines Paares}

F\"ur ein Paar $(X, A)$ mit $A \subseteq X$ gibt es eine lange exakte Sequenz
\[
  \cdots \to H_n(A) \xrightarrow{i_*} H_n(X)
  \xrightarrow{j_*} H_n(X, A)
  \xrightarrow{\partial} H_{n-1}(A) \to \cdots.
\]
Dies ist ein Standardwerkzeug f\"ur die induktive Berechnung von Homologie.

%% ============================================================
\section{Zentrale S\"atze: Mayer-Vietoris und Ausschneidung}
\label{sec:schluessel}
%% ============================================================

\subsection{Die Mayer-Vietoris-Sequenz}

\begin{theorem}[Mayer-Vietoris]
\label{thm:mayer_vietoris}
Sei $X = A \cup B$, wobei $A$ und $B$ offene Teilmengen sind (oder allgemeiner,
wenn die Innengebiete $X$ \"uberdecken). Dann gibt es eine lange exakte Sequenz
\[
  \cdots \to H_n(A \cap B)
  \xrightarrow{\Phi} H_n(A) \oplus H_n(B)
  \xrightarrow{\Psi} H_n(X)
  \xrightarrow{\Delta} H_{n-1}(A \cap B) \to \cdots,
\]
wobei $\Phi(\alpha) = (i_*(\alpha), -j_*(\alpha))$,
$\Psi(\alpha, \beta) = k_*(\alpha) + l_*(\beta)$ und $\Delta$ der
verbindende Homomorphismus ist.
\end{theorem}
\begin{proof}
Man betrachtet die kurze exakte Sequenz von Kettenkomplexen
\[
  0 \to C_*(A \cap B) \xrightarrow{\Phi} C_*(A) \oplus C_*(B)
  \xrightarrow{\Psi} C_*(A+B) \to 0,
\]
wobei $C_*(A+B)$ Ketten bezeichnet, deren Tr\"ager in $A$ oder $B$ liegen.
Die Inklusion $C_*(A+B) \hookrightarrow C_*(X)$ ist eine Quasi-Isomorphie
durch baryzentrischen Unterterteilung (jeder Zyklus in $C_*(X)$ ist
homolog zu einem in $C_*(A+B)$ nach gen\"ugend feiner Unterteilung).
Das Schlangenlemma auf die kurze exakte Sequenz von Kettenkomplexen
angewandt ergibt dann die gew\"unschte lange exakte Sequenz.
\end{proof}

\begin{example}
Zur Berechnung von $H_n(S^n)$: Man schreibt $S^n = U \cup V$, wobei
$U$ und $V$ Halbkugeln sind (beide zusammenziehbar), mit $U \cap V \simeq S^{n-1}$.
Mayer-Vietoris ergibt f\"ur $n \geq 1$:
$H_k(S^n) \cong H_{k-1}(S^{n-1})$ f\"ur $k > 0$ und $H_0(S^n) = \mathbb{Z}$,
woraus das Ergebnis per Induktion folgt.
\end{example}

\subsection{Der Ausschneidungssatz}

\begin{theorem}[Ausschneidungssatz]
\label{thm:ausschneidung}
Sei $Z \subseteq A \subseteq X$ mit $\overline{Z} \subseteq \mathrm{int}(A)$.
Dann induziert die Inklusion $(X \setminus Z, A \setminus Z) \hookrightarrow (X, A)$
Isomorphismen
\[
  H_n(X \setminus Z, A \setminus Z) \xrightarrow{\cong} H_n(X, A)
  \quad \text{f\"ur alle } n.
\]
\end{theorem}
\begin{proof}
Durch baryzentrischen Unterterteilung l\"asst sich jede singul\"are Kette in
$C_n(X, A)$ durch eine Summe von Simplizes darstellen, von denen jeder ganz
in $A$ oder ganz in $X \setminus Z$ liegt. Dies zeigt, dass die Inklusion
$C_*(X \setminus Z, A \setminus Z) \hookrightarrow C_*(X, A)$ eine Surjektion
auf der Homologie induziert. Die Injektivit\"at folgt analog durch Unterteilung
null-homologer Zyklen.
\end{proof}

\subsection{Die Euler-Charakteristik}

\begin{definition}
Die \emph{Euler-Charakteristik} eines endlichen CW-Komplexes $X$ ist
\[
  \chi(X) = \sum_{k=0}^{\dim X} (-1)^k \beta_k,
\]
wobei $\beta_k = \operatorname{rank}(H_k(X;\mathbb{Z}))$ die
\emph{Betti-Zahlen} sind.
\end{definition}

\begin{theorem}[Euler-Poincar\'e-Formel]
F\"ur einen endlichen CW-Komplex mit $c_k$ Zellen der Dimension $k$:
\[
  \chi(X) = \sum_k (-1)^k c_k = \sum_k (-1)^k \beta_k.
\]
\end{theorem}

%% ============================================================
\section{Kohomologie}
\label{sec:kohomologie}
%% ============================================================

\subsection{Kokettenkomplexe}

\begin{definition}
F\"ur eine abelsche Gruppe $G$ ist die \emph{singul\"are Kokettengruppe}
\[
  C^n(X; G) = \mathrm{Hom}(C_n(X;\mathbb{Z}), G).
\]
Der \emph{Korandoperator} $\delta^n\colon C^n \to C^{n+1}$ ist das Transponierte
von $\partial_{n+1}$: $(\delta^n f)(\sigma) = f(\partial_{n+1}\sigma)$.
\end{definition}

Es gilt $\delta^{n+1} \circ \delta^n = 0$, was den Kokettenkomplex
$(C^*(X;G), \delta)$ ergibt.

\begin{definition}
Die \emph{$n$-te Kohomologiegruppe} ist
\[
  H^n(X; G) = \ker \delta^n / \operatorname{im} \delta^{n-1}.
\]
\end{definition}

\subsection{Der universelle Koeffizientensatz}

\begin{theorem}[Universeller Koeffizientensatz f\"ur Kohomologie]
\label{thm:uct}
F\"ur jeden topologischen Raum $X$ und jede abelsche Gruppe $G$ gibt es eine
(nicht nat\"urlich) gespaltene kurze exakte Sequenz
\[
  0 \to \mathrm{Ext}^1_\mathbb{Z}(H_{n-1}(X), G)
  \to H^n(X; G)
  \to \mathrm{Hom}(H_n(X), G)
  \to 0.
\]
\end{theorem}
\begin{proof}
Der Kettenkomplex $C_*(X;\mathbb{Z})$ besteht aus freien abelschen Gruppen.
Anwenden von $\mathrm{Hom}(-,G)$ auf die kurzen exakten Sequenzen
$0 \to Z_n \to C_n \to B_{n-1} \to 0$ und $0 \to B_n \to Z_n \to H_n \to 0$
und Verwendung der langen exakten Sequenzen in Ext (die abbrechen, da $Z_n$
und $B_n$ frei sind) ergibt die gew\"unschte Spaltung.
\end{proof}

\subsection{Die K\"unneth-Formel}

\begin{theorem}[K\"unneth-Formel]
\label{thm:kuenneth}
F\"ur R\"aume $X$ und $Y$ mit freiabelscher Homologie $H_*(X;\mathbb{Z})$:
\[
  H_n(X \times Y;\mathbb{Z})
  \cong \bigoplus_{i+j=n} H_i(X;\mathbb{Z}) \otimes H_j(Y;\mathbb{Z}).
\]
F\"ur Koeffizienten in einem K\"orper $\mathbf{k}$:
\[
  H_n(X \times Y; \mathbf{k})
  \cong \bigoplus_{i+j=n} H_i(X;\mathbf{k}) \otimes_\mathbf{k} H_j(Y;\mathbf{k}).
\]
\end{theorem}
\begin{proof}
Man konstruiert das \emph{Tensorprodukt von Kettenkomplexen}
$(C_*(X) \otimes C_*(Y), \partial^\otimes)$ mit
$\partial^\otimes(\sigma \otimes \tau) = \partial\sigma \otimes \tau
+ (-1)^{|\sigma|} \sigma \otimes \partial\tau$.
Der Eilenberg-Zilber-Satz liefert eine Kett\-en\"aquivalenz
$C_*(X \times Y) \simeq C_*(X) \otimes C_*(Y)$.
Bei freier Homologie $H_*(X)$ vervollständigt die algebraische K\"unneth-Formel
$H_n(A \otimes B) \cong \bigoplus_{i+j=n} H_i(A) \otimes H_j(B)$
(f\"ur freie Kettenkomplexe) den Beweis.
\end{proof}

%% ============================================================
\section{Poincar\'e-Dualit\"at}
\label{sec:poincare}
%% ============================================================

\subsection{Orientierung und Fundamentalklasse}

\begin{definition}
Eine zusammenh\"angende, geschlossene, orientierbare $n$-Mannigfaltigkeit $M$
hat einen eindeutigen Erzeuger $[M] \in H_n(M;\mathbb{Z}) \cong \mathbb{Z}$,
ihre \emph{Fundamentalklasse}.
\end{definition}

\subsection{Cap-Produkt}

Das \emph{Cap-Produkt} $\frown\colon H_k(X;R) \otimes H^l(X;R) \to H_{k-l}(X;R)$
ist der Mechanismus der Poincar\'e-Dualit\"at.

\begin{theorem}[Poincar\'e-Dualit\"at]
\label{thm:poincare_dualitaet}
Sei $M$ eine kompakte, zusammenh\"angende, orientierbare $n$-Mannigfaltigkeit
(ohne Rand) mit Fundamentalklasse $[M] \in H_n(M;\mathbb{Z})$. Dann induziert
das Cap-Produkt mit $[M]$ Isomorphismen
\[
  \mathcal{D}\colon H^k(M;R) \xrightarrow{\cong} H_{n-k}(M;R)
  \quad \text{f\"ur alle } k \text{ und alle Ringe } R,
\]
gegeben durch $\mathcal{D}(\alpha) = [M] \frown \alpha$.
\end{theorem}
\begin{proof}
Der Beweis verl\"auft per Induktion \"uber eine gute \"Uberdeckung. F\"ur einen
offenen Ball $U \cong \mathbb{R}^n$ reduziert sich beides auf die Homologie eines
Punktes, und die Isomorphie ist unmittelbar. F\"ur eine allgemeine Mannigfaltigkeit
\"uberdeckt man $M$ durch euklid-ische Karten und wendet Mayer-Vietoris induktiv
an (das F\"unf-Lemma zeigt, dass das Cap-Produkt bei jedem Schritt ein Isomorphismus
bleibt). Da $M$ kompakt ist, kann die \"Uberdeckung als endlich gew\"ahlt werden,
was die Induktion abschlie\ss{}t.
\end{proof}

\begin{corollary}
F\"ur eine kompakte orientierbare $n$-Mannigfaltigkeit $M$:
$\beta_k(M) = \beta_{n-k}(M)$ f\"ur alle $k$.
Insbesondere gilt $\chi(M) = 0$, falls $n$ ungerade ist.
\end{corollary}

\begin{theorem}[Poincar\'e-Vermutung, Dimension 3 -- Perelman 2003]
\label{thm:poincare_vermutung}
Jede geschlossene, einfach zusammenh\"angende 3-Mannigfaltigkeit ist
hom\"omorph zur $S^3$.
\end{theorem}
\begin{proof}[Historische Anmerkung]
Die Vermutung wurde 1904 von Poincar\'e aufgestellt. Perelman bewies sie
2002--2003 mithilfe des Ricci-Flusses mit Chirurgie von Hamilton. Das Argument
zeigt, dass jede einfach-zusammenh\"angende geschlossene 3-Mannigfaltigkeit
unter dem Ricci-Fluss mit Chirurgie in endlicher Zeit ausgeht (``extinction''),
und die einzige geschlossene 3-Mannigfaltigkeit, die in endlicher Zeit ausgehen
kann, ist (hom\"omorph zu) $S^3$.
\end{proof}

%% ============================================================
\section{CW-Komplexe und zellul\"are Homologie}
\label{sec:cw}
%% ============================================================

\subsection{Definition von CW-Komplexen}

\begin{definition}
Ein \emph{CW-Komplex} ist ein topologischer Raum $X$, der induktiv aufgebaut wird:
\begin{enumerate}[label=(\roman*)]
  \item $X^{-1} = \emptyset$.
  \item $X^n$ entsteht aus $X^{n-1}$ durch Anheften von $n$-Zellen via
    \emph{Anheftungsabbildungen} $\phi_\alpha\colon S^{n-1} \to X^{n-1}$:
    \[
      X^n = X^{n-1} \sqcup_{\bigsqcup_\alpha \phi_\alpha}
        \bigsqcup_\alpha D^n_\alpha.
    \]
  \item $X = \bigcup_n X^n$ mit der schwachen Topologie.
\end{enumerate}
\end{definition}

\subsection{Zellul\"arer Kettenkomplex und Homologie}

\begin{definition}
Die \emph{zellul\"are Kettengruppe} $C_n^{\mathrm{CW}}(X)$ ist die freie abelsche
Gruppe, erzeugt von den $n$-Zellen von $X$. Der zellul\"are Randoperator
$d_n\colon C_n^{\mathrm{CW}} \to C_{n-1}^{\mathrm{CW}}$ ist durch den
\emph{Abbildungsgrad} der Anheftungsabbildungen gegeben:
\[
  d_n(e^n_\alpha) = \sum_\beta d_{\alpha\beta} e^{n-1}_\beta,
\]
wobei $d_{\alpha\beta} \in \mathbb{Z}$ der Grad der Komposition
$S^{n-1} \xrightarrow{\phi_\alpha} X^{n-1} \twoheadrightarrow S^{n-1}_\beta$
ist.
\end{definition}

\begin{theorem}[Zellul\"are Homologie gleich singul\"are Homologie]
$H_n(C_*^{\mathrm{CW}}(X), d) \cong H_n(X;\mathbb{Z})$ f\"ur alle $n$.
\end{theorem}

\begin{example}
CW-Strukturen von Standardr\"aumen:
\begin{center}
\begin{tabular}{lll}
\toprule
Raum & Zellenstruktur & Homologie \\
\midrule
$S^n$ & eine 0-Zelle, eine $n$-Zelle & $H_0 = H_n = \mathbb{Z}$, Rest 0 \\
$\mathbb{R}P^n$ & je eine Zelle pro Dim. $0,\ldots,n$ & abwechselnd $\mathbb{Z}$ und $\mathbb{Z}/2$ \\
$T^2$ & eine 0-, zwei 1-, eine 2-Zelle & $H_0=\mathbb{Z}$, $H_1=\mathbb{Z}^2$, $H_2=\mathbb{Z}$ \\
$\mathbb{R}P^2$ & eine 0-, eine 1-, eine 2-Zelle & $H_0=\mathbb{Z}$, $H_1=\mathbb{Z}/2$, $H_2=0$ \\
\bottomrule
\end{tabular}
\end{center}
\end{example}

\begin{example}[Kompakte Fl\"achen]
Durch CW-Zerlegungen bestimmt man die Homologie aller kompakten Fl\"achen:
\begin{align*}
  &\text{Sph\"are } S^2:\quad H_0=\mathbb{Z},\; H_1=0,\; H_2=\mathbb{Z} \\
  &\text{Torus } T^2:\quad H_0=\mathbb{Z},\; H_1=\mathbb{Z}^2,\; H_2=\mathbb{Z} \\
  &\text{Geschlecht-}g\text{-Fl\"ache } \Sigma_g:\quad H_1=\mathbb{Z}^{2g},\; \chi=2-2g \\
  &\text{RP}^2:\quad H_0=\mathbb{Z},\; H_1=\mathbb{Z}/2,\; H_2=0 \\
  &\text{Kleinsche Flasche } K:\quad H_1=\mathbb{Z}\oplus\mathbb{Z}/2,\; H_2=0
\end{align*}
\end{example}

%% ============================================================
\section{H\"ohere Homotopiegruppen}
\label{sec:hoehere}
%% ============================================================

\subsection{Definition}

\begin{definition}
Die \emph{$n$-te Homotopiegruppe} von $(X, x_0)$ ist
\[
  \pi_n(X, x_0) = [(I^n, \partial I^n), (X, x_0)],
\]
die Menge der Homotopieklassen von Abbildungen
$(I^n, \partial I^n) \to (X, x_0)$ (relativ dem Rand).
\end{definition}

F\"ur $n \geq 2$ ist $\pi_n(X, x_0)$ abelsch. F\"ur $n = 1$ ergibt sich
die Fundamentalgruppe.

\begin{example}
$\pi_n(S^n) \cong \mathbb{Z}$ f\"ur alle $n \geq 1$ (Erzeuger: Identit\"atsabbildung).

$\pi_3(S^2) \cong \mathbb{Z}$ (erzeugt durch die Hopf-Faserung $\eta\colon S^3 \to S^2$).
\end{example}

\begin{remark}
H\"ohere Homotopiegruppen sind im Allgemeinen sehr schwer zu berechnen.
So ist $\pi_k(S^n)$ f\"ur $k > n$ in der Regel nicht-trivial (stabile
Homotopiegruppen der Sph\"aren), und viele davon sind unbekannt.
\end{remark}

\subsection{Der Hurewicz-Isomorphismussatz}

\begin{theorem}[Hurewicz-Isomorphismus]
\label{thm:hurewicz}
Sei $X$ ein $(n-1)$-zusammenh\"angender Raum (d.\,h. $\pi_k(X) = 0$
f\"ur $0 \leq k < n$) mit $n \geq 2$. Dann ist $H_k(X) = 0$ f\"ur $0 < k < n$,
und der \emph{Hurewicz-Homomorphismus}
\[
  h\colon \pi_n(X, x_0) \to H_n(X;\mathbb{Z}),
  \quad [f] \mapsto f_*(\iota_n),
\]
(wobei $\iota_n \in H_n(S^n) \cong \mathbb{Z}$ die Fundamentalklasse ist)
ist ein Isomorphismus.
\end{theorem}
\begin{proof}
Der Beweis verwendet die Serre-Spektralsequenz f\"ur die Weg-Schleifen-Faserung
$\Omega X \to PX \to X$ (wobei $PX$ zusammenziehbar ist). Da $X$ $(n-1)$-zusammenh\"angend
ist, gilt $H_k(\Omega X) = 0$ f\"ur $0 < k < n-1$, und die Spektralsequenz
kollabiert im relevanten Bereich auf den gew\"unschten Isomorphismus.
F\"ur $n=2$ (einfach-zusammenh\"angend) gen\"ugt ein direktes simplizielles Argument:
Jeder 1-Zyklus definiert eine Schleife, die \"uber den Hurewicz-Homomorphismus
einen Isomorphismus ergibt.
\end{proof}

\begin{theorem}[Whiteheads Satz]
Sei $f\colon X \to Y$ eine Abbildung zwischen einfach-zusammenh\"angenden
CW-Komplexen. Induziert $f_*\colon H_n(X;\mathbb{Z}) \xrightarrow{\cong} H_n(Y;\mathbb{Z})$
f\"ur alle $n$, so ist $f$ eine Homotopie\"aquivalenz.
\end{theorem}

\begin{theorem}[Adams -- Hopf-Invariante 1]
\label{thm:adams}
Abbildungen $f\colon S^{2n-1} \to S^n$ der Hopf-Invariante $1$ existieren nur
f\"ur $n = 1, 2, 4, 8$.
\end{theorem}
\begin{proof}[Historische Anmerkung]
Bewiesen von J.\,F.~Adams 1960 mithilfe sekund\"arer Kohomologieoperationen
(Adams-Operationen in der K-Theorie liefern einen eleganteren Beweis).
Das Resultat impliziert, dass $\mathbb{R}^n$ genau f\"ur $n = 1, 2, 4, 8$
eine normierte Divisionsalgebrenstruktur tr\"agt ($\mathbb{R}, \mathbb{C},
\mathbb{H}, \mathbb{O}$).
\end{proof}

%% ============================================================
\section{Anwendungen}
\label{sec:anwendungen}
%% ============================================================

\subsection{Der Brouwer-Fixpunktsatz}

\begin{theorem}[Brouwer-Fixpunktsatz]
\label{thm:brouwer}
Jede stetige Abbildung $f\colon D^n \to D^n$ besitzt einen Fixpunkt.
\end{theorem}
\begin{proof}
Angenommen, $f$ hat keinen Fixpunkt. Definiere eine Retraktion
$r\colon D^n \to S^{n-1}$: F\"ur jedes $x \in D^n$ sei $r(x)$ der Punkt,
in dem der Strahl von $f(x)$ durch $x$ die Sph\"are $S^{n-1}$ trifft.
Diese Abbildung ist stetig und erf\"ullt $r \circ \iota = \mathrm{id}_{S^{n-1}}$,
wobei $\iota\colon S^{n-1} \hookrightarrow D^n$ die Inklusion ist.
Auf der $(n-1)$-ten Homologie ergibt sich:
\[
  \mathbb{Z} \cong H_{n-1}(S^{n-1})
  \xrightarrow{r_*} H_{n-1}(D^n) = 0
  \xrightarrow{\iota_*} H_{n-1}(S^{n-1}) \cong \mathbb{Z},
\]
mit $r_* \circ \iota_* = \mathrm{id}$. Dies erfordert eine Injektion von
$\mathbb{Z}$ in $0$ -- ein Widerspruch.
\end{proof}

\subsection{Der Satz von Borsuk-Ulam}

\begin{theorem}[Borsuk-Ulam]
\label{thm:borsuk_ulam}
F\"ur jede stetige Abbildung $f\colon S^n \to \mathbb{R}^n$ gibt es einen Punkt
$x \in S^n$ mit $f(x) = f(-x)$.
\end{theorem}
\begin{proof}
Angenommen $f(x) \neq f(-x)$ f\"ur alle $x$. Definiere
$g\colon S^n \to S^{n-1}$ durch
$g(x) = (f(x) - f(-x)) / \|f(x) - f(-x)\|$.
Dann gilt $g(-x) = -g(x)$, also ist $g$ eine \emph{ungerade} Abbildung.
Eine ungerade Abbildung $S^n \to S^{n-1}$ induziert beim \"Aquatorschnitt
eine Abbildung ungeraden Grades auf $H_{n-1}$ -- insbesondere nicht die
Nullabbildung. Da jedoch jede Abbildung $S^n \to S^{n-1}$ (f\"ur $n \geq 2$)
auf $H_n(S^n) \neq 0$ trivial wirkt, ergibt sich ein Widerspruch aus der
Lefschetz-Hopf-Theorie.
\end{proof}

\begin{corollary}
Zu jedem Zeitpunkt gibt es auf der Erdoberfl\"ache zwei antipodische Punkte mit
gleicher Temperatur und gleichem Luftdruck.
\end{corollary}

\subsection{Der Igel-Satz (Hairy Ball Theorem)}

\begin{theorem}[Igel-Satz]
\label{thm:igelsatz}
Auf $S^{2n}$ (der geraddimensionalen Sph\"are) existiert kein nirgends
verschwindendes stetiges Tangentialvektorfeld.
\end{theorem}
\begin{proof}
Angenommen, $v\colon S^{2n} \to \mathbb{R}^{2n+1}$ sei ein nirgends
verschwindendes Tangentialvektorfeld (also $v(x) \perp x$ und $v(x) \neq 0$
f\"ur alle $x \in S^{2n}$). Normierung auf $\|v(x)\| = 1$. Definiere eine
Homotopie
\[
  H_t(x) = \cos(\pi t)\, x + \sin(\pi t)\, v(x), \quad t \in [0,1].
\]
Dann ist $H_0 = \mathrm{id}_{S^{2n}}$ und $H_1 = -\mathrm{id}_{S^{2n}}$
(die antipodische Abbildung). Da $H_0$ und $H_1$ homotop sind, haben sie
denselben Abbildungsgrad. Der Grad der Identit\"at ist $+1$, der Grad der
antipodischen Abbildung auf $S^{2n}$ ist $(-1)^{2n+1} = -1$.
Also $1 = -1$ -- Widerspruch.
\end{proof}

\begin{remark}
Auf $S^{2n-1}$ (ungeraddimensionale Sph\"aren) existieren nirgends
verschwindende Vektorfelder: Das Einheitsvektorfeld
$v(x_1,\ldots,x_{2n}) = (-x_2,x_1,\ldots,-x_{2n},x_{2n-1})$
liefert ein solches.
\end{remark}

\subsection{Klassifikation kompakter Fl\"achen}

\begin{theorem}[Klassifikationssatz kompakter Fl\"achen]
Jede kompakte, zusammenh\"angende Fl\"ache ist hom\"omorph zu genau einer der
folgenden:
\begin{enumerate}[label=(\roman*)]
  \item $\Sigma_g$ (orientierbar, Geschlecht $g \geq 0$): $\chi = 2-2g$.
  \item $N_k$ (nicht-orientierbar, $k \geq 1$ Kreuzkappenröhren): $\chi = 2-k$.
\end{enumerate}
Die Euler-Charakteristik und Orientierbarkeit sind vollst\"andige Invarianten.
\end{theorem}

\subsection{Verbindung zur Implementierung}

Das Python-Modul \texttt{algebraic\_topology.py} implementiert die zentralen
Algorithmen dieses Papers:
\begin{itemize}
  \item \texttt{compute\_smith\_normal\_form}: Smith-Normalform für ganzzahlige
    Matrizen -- Grundlage aller Homologieberechnungen.
  \item \texttt{SimplicialComplex}: Simplizialkomplexe mit Randoperatoren,
    Euler-Charakteristik und Betti-Zahlen.
  \item \texttt{SimplicialHomology}: Vollst\"andige Homologieberechnung mit
    Torsionsanteilen.
  \item \texttt{classify\_surface}: Klassifikation kompakter Fl\"achen.
  \item \texttt{van\_kampen\_free\_product}: Seifert-van-Kampen-Satz (Spezialfall).
  \item \texttt{classifying\_space\_demo}: Klassifizierende R\"aume $BG$.
\end{itemize}

%% ============================================================
\section*{Literaturverzeichnis}
%% ============================================================

\begin{thebibliography}{99}

\bibitem{Hatcher}
A.~Hatcher.
\newblock \textit{Algebraic Topology}.
\newblock Cambridge University Press, Cambridge, 2002.
\newblock Frei online: \url{https://pi.math.cornell.edu/~hatcher/AT/ATpage.html}.

\bibitem{Bredon}
G.~E.~Bredon.
\newblock \textit{Topology and Geometry}.
\newblock Graduate Texts in Mathematics 139, Springer, New York, 1993.

\bibitem{May}
J.~P.~May.
\newblock \textit{A Concise Course in Algebraic Topology}.
\newblock University of Chicago Press, Chicago, 1999.

\bibitem{Munkres}
J.~R.~Munkres.
\newblock \textit{Elements of Algebraic Topology}.
\newblock Addison-Wesley, Menlo Park, CA, 1984.

\bibitem{Spanier}
E.~H.~Spanier.
\newblock \textit{Algebraische Topologie} (Algebraic Topology).
\newblock Springer, New York, 1981.

\bibitem{Adams}
J.~F.~Adams.
\newblock On the non-existence of elements of Hopf invariant one.
\newblock \textit{Annals of Mathematics}, 72:20--104, 1960.

\bibitem{Perelman1}
G.~Perelman.
\newblock The entropy formula for the Ricci flow and its geometric applications.
\newblock \textit{arXiv:math/0211159}, 2002.

\bibitem{Perelman2}
G.~Perelman.
\newblock Ricci flow with surgery on three-manifolds.
\newblock \textit{arXiv:math/0303109}, 2003.

\bibitem{Milnor}
J.~Milnor.
\newblock \textit{Morse Theory}.
\newblock Princeton University Press, Princeton, 1963.

\end{thebibliography}

\end{document}
